\section{Conclusion}

This paper advocates for processors to incorporate microarchitectural isolation primitives \emph{and} mechanisms for controlled speculation, enabling different execution modes.
These modes can restrict what is exposed to an attacker, 
effectively balancing performance and program-analysis complexity.
We introduce two mechanisms to securely share memory between an enclave and an untrusted OS in an out-of-order processor.
We show through security analysis and performance evaluation, that these two modes are complementary, achieving speculative non-interference with a reasonable performance impact, while requiring minimal code annotation and simple program analysis doable by hand.
Our prototype, Citadel, is a multicore processor running on an FPGA, booting untrusted Linux, and supporting comprehensive enclave capabilities, such as shared memory, and remote attestation.
To our knowledge, Citadel is the first end-to-end enclave platform to run secure applications, such as cryptographic libraries or small private inference workloads, on a speculative out-of-order multicore processor.